\documentclass[10pt,a4paper]{amsart}
\usepackage[margin=19mm]{geometry}
\usepackage{amsmath,amssymb,mathtools}
\usepackage[scr=rsfs,cal=euler]{mathalfa}
\usepackage{xfrac}
\usepackage[unicode,hidelinks,bookmarks=false]{hyperref}
\newcommand{\ie}{i.e.\ }
\newcommand{\cf}{cf.\ }
\newcommand{\resp}{respectively\ }
\newcommand{\Subsection}{\subsection}
\providecommand{\nobreakdash}{\nobreak\hbox{-}}
\setlength{\emergencystretch}{2em}
\multlinegap0pt
\usepackage{bm}
\usepackage{bbm}
\usepackage{dsfont}
\usepackage{xspace}
\usepackage{cancel}
\usepackage{mathtools}
\usepackage{amsthm}
\makeatletter
\@ifundefined{theorem}{\newtheorem{theorem}{Theorem}}{}
\@ifundefined{lemma}{\newtheorem{lemma}[theorem]{Lemma}}{}
\makeatother
\newcommand\CP{\mathbb{CP}}
\newcommand\C{\mathbb{C}}
\newcommand\N{\mathbb{N}}
\newcommand\U{\mathbb{U}}
\newcommand\di{\partial}
\newcommand\wt{\widetilde}
\title[Every projective Oka manifold is elliptic]{Every projective Oka manifold is elliptic}
\author{Franc Forstneric, Finnur Larusson}
\date{Source: arXiv:2502.20028v6 (CC BY 4.0)}
\begin{document}
\maketitle

\noindent\emph{Source and licence.} Every projective Oka manifold is elliptic. Version arXiv:2502.20028v6 (CC BY 4.0), distributed under the Creative Commons Attribution 4.0 International licence, as recorded in the arXiv licence metadata for this exact version and shown on the abstract page https://arxiv.org/abs/2502.20028v6. Full rights notice: licenses/arxiv-2502.20028-CC-BY-4.0.txt.\par\smallskip

\noindent\textit{Editorial scope.} Counted unit: the Lemma whose source label is \texttt{lem:V} in the article (source0.tex lines 544--552, source label \texttt{lem:V}), delivered together with the article's own complete original proof of it (source0.tex lines 554--657, one \texttt{proof} environment of 104 lines). No mathematics is written, repaired or re-derived: statement and proof are the article's own text line for line. Required context retained uncounted: eq:V (lines 526--528). Every definition, display and label of those lines is retained unchanged. Omitted from this package and not redistributed: 1--525, 529--543, 553--553, 658--1303. Nothing inside a retained excerpt is omitted or altered: every retained body is delivered exactly as the article's lines are written, so no omission receipt of any kind accompanies this package. Citations: the retained excerpts carry no citation command at all, so no bibliography entry of the article is transcribed and no reference is redistributed. Reference adaptation: none was needed and none was made; every reference inside the delivered unit resolves inside it, so no unresolved reference or citation is produced. Printed slips kept literal: no printed word, symbol, hypothesis or number of the counted statement or of its proof was altered. Layout and class: the article's own class is \texttt{amsart} with packages including babel, inputenc, amsfonts, amssymb, amscd, amstext, mathrsfs, verbatim, enumerate; the wrapper is the lane's pinned \texttt{amsart} editorial wrapper at 10pt on A4, which restates the theorem-like environments the retained text needs and adds the article's own macro definitions that the retained text needs (\texttt{\textbackslash C}, \texttt{\textbackslash CP}, \texttt{\textbackslash N}, \texttt{\textbackslash U}, \texttt{\textbackslash di}, \texttt{\textbackslash wt}), with extra wrapper packages (bm, bbm, dsfont, xspace, cancel, mathtools). The delivered numbering is the unit's own. Assets: no retained span contains a figure environment, a graphics-inclusion command, a PDF-page inclusion, a table or a TikZ picture; no image file is copied into this package, the package declares no asset dependency and no image file is redistributed with it. Licence: version arXiv:2502.20028v6 of this article is distributed under the Creative Commons Attribution 4.0 International licence, as recorded in the arXiv licence metadata for this exact version and shown on the abstract page https://arxiv.org/abs/2502.20028v6. A full rights notice is delivered at licenses/arxiv-2502.20028-CC-BY-4.0.txt. Characters: every retained excerpt is pure ASCII, so the delivered engine, XeLaTeX with its default Latin Modern text font, reports no missing character.
\begin{equation}\label{eq:V}
	V(x,t) = \sum_{i=1}^n \sum_{j=1}^m t_j \, V_{i,j}(x) \di_{x_i},
\end{equation}

\begin{lemma}\label{lem:V}
Let $k_0$ denote the maximum of the degrees of the polynomials $V_{i,j}(x)$.
For every $k\ge k_0$, the vector field $V$
\eqref{eq:V} extends to an algebraic vector field on the total space
of the vector bundle $E=(\CP^n \times \C^m)\otimes \U^k= m\U^k$ on 
$\CP^n$ (the direct sum of $m$ copies of the $k$-th tensor power of 
$\U$), which vanishes on the zero section
$E_0$ of $E$ and on $E|_{\Lambda_0}$. 
\end{lemma}

\begin{proof}
For every $\alpha=0,1,\ldots,n$ we have a vector bundle trivialisation 
$\theta_\alpha:E|U_\alpha\stackrel{\cong}{\to} U_\alpha\times \C^m$ 
with transition maps 
$\theta_{\alpha,\beta}=\theta_\alpha\circ \theta_\beta^{-1}$ on 
$(U_\alpha\cap U_\beta) \times \C^m$ given by 
\begin{equation}\label{eq:transition}
	\theta_{\alpha,\beta}([z],t) = \bigl([z],(z_\alpha/z_\beta)^k t\bigr),\quad 
	t\in\C^m, \ 0\le \alpha,\beta\le n.
\end{equation}
In particular, $\theta_{\alpha,0}([z],t)=\bigl([z],(z_\alpha/z_0)^k t\bigr)$.
We analyse the behaviour of the vector field $V$ \eqref{eq:V} 
near the hyperplane
$\Lambda_0\setminus \Lambda_\alpha$ for $\alpha=1,\ldots,n$. 
It suffices to consider the case $\alpha=1$ since the same argument 
will apply to every $\alpha$. 
For simplicity we first make the calculation in the case 
when $m=1$, so $E=\U^k$. The calculation is made in two steps. 
In the first step, we express $V$ in the fibre coordinate $t'$ on the 
line bundle chart $E|{U_1}\cong \C^n\times \C$ over the domain
$U_0\cap U_1 =\{x=(x_1,\ldots,x_n)\in \C^n: x_1\ne 0\}$.
The transition map $\theta_{1,0}$ is given by
$\theta_{1,0}(x,t)=\bigl(x,x_1^{k} t\bigr)$, so $t'=x_1^{k} t$. 
Its differential has the block form
\begin{equation}\label{eq:Dtheta}
	D\theta_{1,0}(x,t)=
	\begin{pmatrix}
	I_{n} & 0 \\
	B      & x_1^{k} 
	\end{pmatrix}
\end{equation}
where $I_n$ is the $n\times n$ identity matrix and
$B=(kx_1^{k-1} t,0,\ldots,0)$. It follows that the vector field
$V'=D\theta_{1,0} \, \cdotp V$ equals
\begin{align*} 
	V' &= t\sum_{i=1}^n V_i(x)\di_{x_i} + k t^2 x_1^{k-1}V_1(x) \di_{t'} \\
	   & = t' x_1^{-k}\sum_{i=1}^n V_i(x)\di_{x_i} + 
		k (t')^2 x_1^{-k-1} V_1(x) \di_{t'},
\end{align*} 
where we used that $t=x_1^{-k}t'$.
In the second step, we express the vector field $V'$ in the standard affine
coordinates $x'=(x'_1,x'_2,\ldots,x'_n)$ on $U_1$. Note that
\begin{align*}
	x'_1 &= \frac{z_0}{z_1} = \frac{1}{x_1},\\
	x'_i  &= \frac{z_i}{z_1} = \frac{x_i}{x_1},\quad i=2,\ldots,n. 	
\end{align*}
Write $x'=\psi(x)$ and $(x',t')=\tilde \psi(x,t')=(\psi(x),t')$. 
We have that 
\[
	D\psi(x)=
	\begin{pmatrix}
	-\frac{1}{x_1^2} & 0 & 0 & \cdots & 0 \\
	-\frac{x_2}{x_1^2} & \frac{1}{x_1} & 0 & \cdots & 0 \\
	\cdots & \cdots & \cdots & \cdots & \cdots  \\
	-\frac{x_n}{x_1^2} & 0 & 0 & \cdots & \frac{1}{x_1}
	\end{pmatrix}.
\]
Hence, the vector field $\wt V=D\tilde\psi\,\cdotp V'$ equals
\begin{align*}
	\wt V(x,t') &= -\frac{1}{x_1^{k+2}} V_1(x) t' \di_{x'_1} 
 	+\sum_{i=2}^n \left[ - \frac{x_i}{x_1^{k+2}} V_1(x) 
	+ \frac{1}{x_1^{k+1}} V_i(x)\right] t' \di_{x'_i} \\
	&\quad + \frac{k}{x_1^{k+1}} V_1(x) (t')^2 \di_{t'}. 
\end{align*}	
Note that $x=\psi^{-1}(x')=(1/x'_1,x'_2/x'_1,\cdots,x'_n/x'_1)$.
Inserting this in the above expression gives  
\begin{align*}\label{eq:Vtilde}
	\wt V(x',t') &= -(x'_1)^{k+2} V_1(\psi^{-1}(x')) \, t' \di_{x'_1} \\ 
	&\quad	+\sum_{i=2}^n (x'_1)^{k+1}  \left[ - x'_i V_1(\psi^{-1}(x')) 
		+ V_i(\psi^{-1}(x')) \right]  t' \di_{x'_i} 	\\
	&\quad + k (x'_1)^{k+1} V_1(\psi^{-1}(x')) (t')^2 \di_{t'}.
\end{align*}
Note that the affine hyperplane $\{z_0=0,\ z_1\ne 0\}$ 
corresponds to $\{x'_1=0\}$. 
Since $\psi^{-1}$ is a fractional linear map with a simple
pole along $x'_1=0$, the functions $V_i(\psi^{-1}(x'))$ are rational
in $x'$ with a pole of degree at most $k_0$ along $x'_1=0$ 
and no other singularities. It follows from the above expression 
that for $k\ge k_0$ the vector field $\wt V(x',t')$ is polynomial in $(x',t')$
and it vanishes on $\{x'_1=0\}\cup\{t'=0\}$.

In the general case when $V$ is given by \eqref{eq:V} and $m\in\N$
is arbitrary, the calculation is similar. Using the fibre
variables $t=(t_1,\ldots,t_m)$ on $E|U_0$ and
$t'=(t'_1,\ldots,t'_m)=x_1^k t$ on $E|U_1$, 
the differential $D\theta_{0,1}$ has the block form~\eqref{eq:Dtheta},
where $B$ is now an $m\times n$ matrix and the lower right 
entry is replaced by $x_1^{k}I_m$. The vector field 
$V'=D\theta_{1,0} \, \cdotp V$ equals
\begin{align*}%\label{eq:Vprime}
	V' &= \sum_{i=1}^n \sum_{j=1}^m t_j\, V_{i,j}(x) \di_{x_i} 
	  	 + k x_1^{k-1} \sum_{j,l=1}^m t_jt_l V_{1,j}(x) \di_{t'_l} \\
	   &= x_1^{-k} \sum_{i=1}^n \sum_{j=1}^m t'_j\, V_{i,j}(x) \di_{x_i} 
	   	 + k x_1^{-k-1} \sum_{j,l=1}^m t'_j t'_l V_{1,j}(x) \di_{t'_l}.
\end{align*}
The second step, expressing $V'$ in the affine coordinates $x'$
on $U_1$, is the same as before, and we leave 
the details to the reader. The new vector field 
$\wt V=D\tilde\psi\,\cdotp V'$ 
is polynomial in $(x',t')$, of second order in $t'$, and   
it vanishes on $\{x'_1=0\}\cup\{t'=0\}$.
Since this argument holds on every chart $U_\alpha$ for 
$\alpha=1,\ldots,n$, the lemma is proved.
\end{proof}

\end{document}
